\documentclass[12pt]{article}
\usepackage[utf8]{inputenc}

\usepackage{../mystyle}

\title{Intersection Theory}
\author{Joseph Sullivan}
\date{November 2024}

\tikzset{
    wb/.style={
        ->,preaction={draw, line width=5pt, white}
    }
}

\DeclareMathOperator{\Kom}{Kom}
\DeclareMathOperator{\Cyl}{Cyl}
\DeclareMathOperator{\Stab}{Stab}
\DeclareMathOperator{\Slice}{Slice}
\DeclareMathOperator{\chern}{ch}
\DeclareMathOperator{\todd}{td}

\DeclareMathOperator{\Rat}{Rat}

\newcommand{\cQ}{\mathcal{Q}}
\newcommand{\cZ}{\mathcal{Z}}

\newcommand{\rddots}{\text{\reflectbox{$\ddots$}}}


\begin{document}

\maketitle

\section{Intersections}

\subsection{Motivation}

We recall that in differential topology, we have some intersection theory for submanifolds---given $M,N\hookrightarrow X$ closed oriented submanifolds of dimensions $m,n$ inside a compact oriented manifold $X$ of dimension $d$, we define the intersection of $M,N$ as follows: (1) we homotop the inclusions of $M,N$ until they are transverse, (2) we intersect in $X$ (equivalently pull $N$ back along the inclusion of $M$ and pushforward).

This answer for $M\cap N$ is well-defined up to (co)homology classes. That is, the inclusions $M,N\hookrightarrow X$ give us the cohomology class
\[[M]\in H_m(M) \to H_m(X) \to H^{d-m}(X)\]
and likewise $[N]\in H^{d-n}$. The cohomology classes of $[M],[N]\in H^*(X)$ will be equal to the cohomology classes of homotop'd of inclusions that make $M,N$ transverse. Then, no matter which homotopy we chose to make $M,N$ transverse, we get the same cohomology class $[M\cap N] \in H^{2d-m-n}(X)$.

In fact, it is a theorem (see \cite{cup_products}) that
\[[M]\smile [N] = [M \cap N],\]
so the cohomology ring is a natural setting to study intersection.

Another potential setting for intersection theory is the relative bordism ring {\color{red} (seems really cool!!!)}. This has the advantage of every class really being represented by a smooth manifold (by definition), rather than it sometimes being a theorem for singular cohomology.

We want to define an analogous notion of intersection for subvarieties of smooth quasi-projective variety. 

\subsection{Chow Ring}

\begin{defn}
    Let $X$ be a variety (or more generally a scheme).
    \begin{itemize}
        \item The \textbf{group of $k$-cycles}, denoted $Z_k(X)$, is the free abelian group generated by the set of $k$-dimensional irreducible subvarieties (reduced irreducible subschemes) of $X$.
        \item The \textbf{group of cycles} is the direct sum $Z(X) = \bigoplus_k Z_k(X)$.
        \item A cycle $Z=\sum n_i Y_i$ is called \textbf{effective} if all $n_i\geq 0$.
        \item On a pure n-dimensional scheme, $(n-1)$-cycles are called \textbf{Weil divisors}.
        \item To a closed subscheme $Y\subseteq X$, associate a cycle by taking the irreducible components $Y_1,\dots,Y_s$ of the reduced scheme $Y_{\mathrm{red}}$, and summing
        \[\<Y\> = \sum_i \ell_{\cO_{X,Y_i}}(\cO{Y,Y_i}) Y_i.\]
        where $\ell_i \defeq \operatorname{mult}_{Y_i}(Y) = \ell_{\cO_{X,Y_i}}(\cO{Y,Y_i})$ is called the \textbf{multiplicity} of $Y$ along $Y_i$.
    \end{itemize}
\end{defn}

\begin{rmk}
    Note that associating a subscheme $Y\subseteq X$ to a cycle forgets about embedded primes, since we really sum over irreducible components, not associated primes. This is how it should be---for an example of the dangerous alternative, take $X=\A^2$ and
    \[Y = \Spec k[x,y]/((y)\cap (x,y)^2).\]
    Then,
    \[\ell_{k[x,y]_{(x,y)}} (k[x,y]/((y)\cap (x,y)^2))_{(x,y)} = \infty.\]
    We can see this from an infinite chain of submodules
    \[\<1\> \supsetneq \<\overline{x}\> \supsetneq \<\overline{x}^2\> \supsetneq \cdots.\]
\end{rmk}

Just as in the differential topology case, we do our intersection theory by looking at the free group of singular chains/submanifolds, and then putting an equivalence relation to make intersection well-defined. Unlike singular homology, we don't have to take subobjects first---we just defined the algebraic cycles rather than having to first define chains.

Now, there are several possible equivalence relations on algebraic cycles, the finest of which is \textit{rational equivalence}.

\begin{defn}
    Let $\Rat(X)\subseteq Z(X)$ be the subgroup generated by the differences of the form
    \[\<\Phi \cap (\{t_0\}\times X)\> - \<\Phi \cap (\{t_1\}\times X)\>\]
    where $t_0,t_1\in \P^1$ and $\Phi$ is an irreducible subvariety of $\P^1\times X$ not contained in any fiber $\{t\}\times X$.

    We should think of elements of $\Rat(X)$ like cobordisms (in $X$). Maybe more familiarly, they are analogous to homology boundaries. We say that cycles $A,B$ are \textbf{rationally equivalent} if $B-A\in \Rat(X)$.
\end{defn}

\begin{defn}
    The \textbf{Chow group} of $X$ is the quotient
    \[A(X) = Z(X)/\Rat(X).\]
\end{defn}

\begin{prop}
    The Chow group of $X$ is graded by dimension, i.e.,
    \[A(X) = \bigoplus A_k(X)\]
    with $A_k(X)$ the group of rational equivalence classes of $k$-cycles.
\end{prop}

\begin{proof}
    We show for a triple $(\Phi,t_0,t_1)$ giving a generator of $\Rat(X)$, that the two cross sections are of the same dimension.

    The function $t-t_0$ on $\Phi \cap \A^1\otimes X$ (for $\A^1 \ni t_0$) is a nonzerodivisor, and its zero locus is exactly $\Phi \cap (\{t_0\}\times X)$. So by Krull's principle ideal theorem, the $t_0$ cross section is codimension 1 in $\Phi$.

    Similarly, the $t_1$ cross section of $\Phi$ is codimension 1, so both are of the same dimension.
\end{proof}

\begin{rmk}
    When $X$ is equidimensional, we can define codimension of cycles, and we will denote
    \[A^c(X) = A_{\dim X - c}(X).\]
\end{rmk}

Now that we have the Chow ring, the analog of the singular cohomology ring, we need to define cup product/transverse intersection. We follow along the lines of transverse intersection.

\begin{defn}
    Let $A,B\subseteq X$ be irreducible subvarieties. If $p$ is a smooth point for $A,B,X$, then we say $A$ and $B$ \textbf{intersect transversely} at $p$ if
    \[T_p A + T_p B = T_p X.\]

    We say subvarieties $A,B\subseteq X$ are \textbf{generically transverse} (or they intersect \textbf{generically transversely}) if they meet transversely at a general point of each component $C$ of $A\cap B$.
\end{defn}

{\color{red} I wish we would say ``generally transverse'' instead of ``generically transverse'', because it doesn't seem to be about generic points, but rather this is asserting something happens on a dense Zariski open set.}

\begin{eg}
    Suppose $A,B$ have complementary dimensions in $X$. Then, $A$ and $B$ are generically transverse if and only if they are transverse everywhere. Clearly transverse implies generically transverse.
    
    Conversely, assume $A,B$ are generically transverse. Without loss of generality assume $A,B$ are equidimensional. First, generic transversality implies every component of $A\cap B$ is dimension 0. Otherwise, we take the connected component $C$ of the largest (nonzero) dimension, and by generic transversality there exists a point $P\in C$ where $A,B$ meet transversely. In particular, $A,B,C$ are all smooth at $P$, so
    \[\dim A = \dim T_P A, \qquad \dim B = \dim T_P B, \qquad \dim C = \dim (T_P A\cap T_P B)\]
    and then transversality at $P$ gives
    \[\codim A + \codim B = \codim T_P A + \codim T_P B = \codim(T_P A \cap T_P B) = \codim C,\]
    which implies $\codim C = \dim X$, so $\dim C = 0$.

    Now that $C$ is $0$-dimensional, it is finitely many points, and $A,B$ meet transversely at each of them because they meet transversely at a Zariski dense open set of each connected component of $C$, which is each point.
\end{eg}

\begin{eg}
    If $\codim A + \codim B > \dim X$, then $A,B$ are generically transverse if and only if they are disjoint, since at an intersection point, there is no hope for the sum of tangent spaces to span all of $X$.
\end{eg}

\begin{thmdef}
    If $X$ is a smooth quasi-projective variety, then there is a unique product structure on $A(X)$ satisfying the following condition: If two irreducible subvarities $A,B\subseteq X$ are generically transverse, then
    \[[A][B]=[A\cap B].\]
    This structure makes
    \[A(X) = \bigoplus_{c=0}^{\dim X} A^c(X)\]
    into an associative, commutative ring, graded by codimension, called the Chow ring of $X$.
\end{thmdef}

\begin{lem}
    
\end{lem}

\section{Chern Classes}

In pursuit of classifying vector bundles on a smooth $n$-dimensional quasi-projective variety $X$ over $\C$, we should put down some invariants. For a vector bundle $E$, rather than give numerical invariants, we will give an invariant that lives in the cohomology ring $H^*(X;\C)$, called the \textbf{total Chern class}, denoted $c(E)$. It will only have nonzero components in even degrees. Usually, we will write
\[\bigoplus_{i=1}^{n} H^{2i}(X;\C) \cdot t^i\]
so putting a formal symbol $t^i$ to indicate the degree. In this case, we will instead call the total Chern class the \textbf{Chern polynomial}, and we will denote it by $c_t(E)$.

Writing $\pi_i:H^* \to H^{2i}$ the projection map, we will denote the \textbf{chern classes} $c_i(E) = \pi_i(c_t(E))$.

\begin{prop}
    There exists a unique well-defined assignment, where $E$ is mapped to $c_t(E) \in H^*(X;\C)$ satisfying
    \begin{enumerate}[(i)]
        \item $c(\cO_X(D)) = 1+[D]t$
        \item for $f:X'\to X$, $E$ a vector bundle on $X$, then
        \[c_i(f^* E) = f^*c_i(E).\]
        \item (Whitney sum formula) Given a SES $0\to E' \to E \to E'' \to 0$ of vector bundles, $c_t(E)=c_t(E') c_t(E'')$.
    \end{enumerate}
\end{prop}

A few remarks are in order:
\begin{itemize}
    \item Maybe the moral of Poincare duality on a compact manifold (e.g. a smooth projective variety over $\C$) is that cohomology classes are geometric objects (dual to cycles), and the cup product is intersection. So, the total Chern class turns a vector bundle $E$ into a geometric object inside $X$, and there will be some relation between vector bundle algebra and intersection theory in $X$.
    \item In (i), we mean $[D]\in H^2(X;\Z)$ is the Poincare dual of the homology class of the cycle $[D] H_{2n-2}(X;\Z)$. For example, when $D$ is smooth, it is a $n-1$ complex manifold, so in particular a $2n-2$ orientable manifold with an orientation class $[D]\in H_{2n-2}(D;\Z)$ which under the inclusion $i:D\to X$ we can pushforward $i_*[D] \in H_{2n-2}(X;\Z)$. When not smooth, $D$ is at least a CW complex, and (apparently) its top homology is $\Z$ and has a canonical generator, again given by the choice of $\sqrt{-1}\in \R[x]/(x^2+1)$ as orientation is given in the smooth case.
    \item As a result of (iii), if we have a filtration
    \[0=E_0 \leq E_1 \leq \cdots \leq E_n=E,\]
    of vector subbundles with successive quotients $F_i = E_i/E_{i-1}$ still vector bundles, then we have short exact sequences
    \[0\to E_{i-1} \to E_i \to F_i \to 0,\]
    which give $c_t(E_i) = c_t(F_i)c_t(E_{i-1})$, so by induction,
    \[c_t(E) = \prod_{i=1}^n c_t(F_i).\]
    \item With the previous remark, the \textbf{splitting principle} is the idea/suggestion that to deduce information about Chern polynomials, one should find a morphism $f:X'\to X$ where $f^*E$ (for each $E$ in a finite collection) has a filtration with associated graded pieces which are line bundles. One should also choose $f$ such that $f^*: H^*(X) \to H^*(X')$ is injective (and this turns out to be possible!), so if we find an equation on $c(f^*E_1)=f^*c(E_1),\dots,c(f^*E_n)=f^*c(E_n)$, we get an equation on on $c(E_1),\dots,c(E_n)$. By this line of reasoning, we can assume that all vector bundles filter by line bundles whenever we want to prove an equation. Therefore, by the previous remark, if we assume $E$ is filtered by $\cO_X(D_i)$, we can write
    \[c_t(E) = \prod_{i=1}^r (1+D_it),\]
    and we will refer to the $D_i$ as the \textbf{Chern roots} of $E$.
    \item When we assume $c_t(E) = \prod_{i=1}^r (1+D_i t)$, we can attain the Chern classes $c_i(E)$ by distributing the product, which exactly says that the Chern classes are the elementary symmetric polynomials in the Chern roots.
    \item The Chern classes aren't quite well-defined on the Grothendieck group, despite interacting with short exact sequences. The problem is that the Chern polynomial turns addition into multiplication. We will modify the Chern polynomial to get the \textit{Chern character}, which \textit{will} be well-defined on the Grothendieck group.
    \item Though the entire Chern class $c$ isn't additive on exact sequences, $c_1$ individually is additive, since comparing coefficients on
    \begin{align*}
        c_t(E) &= c_t(E') c_t(E'')\\
        1 + c_1(E)t + \cdots &= (1+c_1(E')t + \cdots)(1+c_1(E'')t + \cdots)\\
        &= 1 + (c_1(E') + c_1(E''))t + \cdots,
    \end{align*}
    we get $c_1(E) = c_1(E') + c_1(E'')$. 
\end{itemize}

\begin{eg}
    Let's find some formulas for computing the Chern classes of a tensor product. Let $E,E'$ be vector bundles, and by the splitting principle, assume
    \[E = \bigoplus_{i=1}^r \cO_X(D_i), \qquad E' = \bigoplus_{j=1}^{r'} \cO_X(D_j'),\]
    (this isn't strictly correct, we should instead assume only that we have filtrations, but its more inconvenient to write), so
    \[E\otimes E' = \bigoplus_{i,j} \cO_X(D_i + D_j).\]
    Now, we take the Chern polynomials of $E$, $E'$, and $E\otimes E'$. Starting with $E$,
    \begin{align*}
        c_t(E) &= \prod_{i=1}^r (1+D_i t)\\
        &= 1 + \left(\sum_{i_1} D_{i_1}\right) t + \left(\sum_{i_1<i_2} D_{i_1}\cdot D_{i_2}\right)t^2 + \cdots + (D_1\cdots D_r) t^n,
    \end{align*}
    and analogously for $E'$. Notice that in $c_2(E)$, we don't get terms $D_i^2$. Rather, to get these terms we look at $c_1(E)^2$. In fact,
    \[c_1(E)^2 - 2c_2(E) = \sum_i D_i^2.\]
    
    Now, for $E\otimes E'$, we get
    \begin{align*}
        c_t(E\otimes E') &= \prod_{i,j} (1+(D_i + D'_j)t).
    \end{align*}
    We start computing the coefficients of the Chern polynomial, which becomes a combinatorics/elementary symmetric polynomial exercise. To distribute out this product to get the degree $k$ terms, we will form a product of $rr'-k$ copies of $1$, and $k$ terms of the form $(D_i+D'_j)t$. So, the degree $k$ coefficient of the Chern polynomial is the sum of all $k$-fold products of $D_i + D'_j$.
    
    To help think about this, we write down a table.
    \begin{center}
    \begin{tabular}{c|c c c c c}
        $r'$   & $D_1 + D'_{r'}$ & $D_2 +D'_{r'}$ & $\cdots$ & $D_{r-1} +D'_{r'}$ & $D_r + D'_{r'}$ \\
        $r'-1$ & $D_1 + D'_{r'-1}$ & $D_2 +D'_{r'-1}$ & $\cdots$ & $D_{r-1} +D'_{r'-1}$ & $D_r + D'_{r'-1}$ \\
        $\vdots$ & $\vdots$ & $\vdots$ & $\rddots$ & $\vdots$ & $\vdots$ \\
        $2$ & $D_1 + D'_2$ & $D_2 +D'_2$ & $\cdots$ & $D_{r-1} +D'_2$ & $D_r + D'_2$\\
        $1$ & $D_1 + D'_1$ & $D_2 + D'_1$ & $\cdots$ & $D_{r-1} + D'_1$ & $D_r + D'_1$\\\hline
            & $1$ & $2$ & $\cdots$ & $r-1$ & $r$
    \end{tabular}
    \end{center}
    In degree 0 we get $1$, the sum of the only 0-fold product. In degree 1, we sum the 1-fold products to get
    \begin{align*}
        \sum_{i,j}(D_i + D_j') &= r'\sum_{i_1} D_i + r\sum_{j_1} D'_{j_1}\\
        &= r'c_1(E) + rc_1(E').
    \end{align*}
    In degree 2, we get the sum of all 2-fold products of $D_i + D_j'$. A single product looks like
    \[(D_{i_1} + D'_{j_1})(D_{i_2} + D'_{j_2}) = D_{i_1}D_{i_2} + D_{i_1}D'_{j_2} + D'_{j_1}D_{i_2} + D'_{j_1}D'_{j_2}.\]
    Now, we count terms. For $i_1=i_2=i$, we get $D_i^2$ occurring $\binom{r'}{2}$ times. For $i_1\neq i_2$, $D_{i_1}D_{i_2}$ occurs $(r')^2$ times. The term $D_{i_1}D'_{j_2}$ occurs $(r-1)(r'-1) + r-1 + r'-1 = rr'-1$ times. So,
    \begin{align*}
        c_2(E\otimes E') &= \binom{r'}{2} \sum_{i=1}^n D_i^2 + (r')^2\sum_{i_1\neq i_2} D_{i_1} D_{i_2} + \binom{r}{2} \sum_{j=1}^n D_j^2 + r^2 \sum_{j_1 \neq j_2} D'_{j_1}D'_{j_2} + (rr'-1)\sum_{i,j} D_iD'_j\\
        &= \binom{r'}{2}(c_1(E)^2-2c_2(E)) + (r')^2 c_2(E)\\
        &\qquad\qquad\qquad + \binom{r}{2}(c_1(E')^2-2c_2(E')) + r^2 c_2(E') + (rr'-1) c_1(E)c_1(E')\\
        &= \binom{r'}{2} c_1(E)^2 + \binom{r}{2} c_1(E')^2 + r'c_2(E) + rc_2(E') + (rr'-1)c_1(E)c_1(E').
    \end{align*}
    If $E'=L$ is a line bundle (so $r'=1$), then
    \[c_2(E\otimes L) = \binom{r}{2}c_1(L)^2 + c_2(E) + (r-1)c_1(E)c_1(L).\]
    One can imagine that $c_3(E\otimes E')$ is even more complicated.
\end{eg}

\begin{defn}
    For a vector bundle $E$ over $X$, there is an element $\chern(E)\in H^*(X;\Q)$ called the \textbf{Chern character}, defined as follows. By the splitting principle, assume $E$ has Chern roots $a_i=[D_i]$, and so
    \[c_t(E)=\prod_{i=1}^r (1+a_i t).\]
    Then, define the Chern character of the line bundle $E$ by
    \begin{align*}
        \chern(E) &= \sum_{i=1}^r e^{a_i t}\\
        &= \sum_{i=1}^r (1 + a_i t + \frac{1}{2!} a_i^2 t^2 + \cdots + \frac{1}{n!} a_i^n t^n)
    \end{align*}
    where $\dim X = n$, so higher intersections are 0 anyways.

    As with the Chern classes, denote $\chern_i(E)$ the projection of $\chern(E)$ in $H^{2i}(X;\Q)$.
\end{defn}

\begin{rmk}
    Unlike the Chern polynomial, where the constant coefficient $c_0(E)$ is always $1$, we can see from the definition that $\chern_0(E) = r$, the rank of $E$.
    
    For the remaining coefficients of $\chern(E)$,  we notice that the expression for the Chern character is symmetric in the Chern roots, so the coefficients are combinations of the elementary symmetric polynomials in the Chern roots, which are exactly the $c_i(E)$. That is, $\chern_i(E)$ will be some $\Q$-combination of the $c_i(E)$. In some small degrees, we can compute
    \[\chern = r + c_1 t + \frac{1}{2}(c_1^2 - 2 c_2)t^2 + \frac{1}{6}(c_1^3 -3c_1c_2 + c_3)t^3 + \cdots.\]
\end{rmk}


\begin{prop}
    The Chern character is well-defined on the Grothendieck group $K(X)$, i.e., for a short exact sequence
    \[0\to E' \to E \to E'' \to 0,\]
    we get $\chern(E) = \chern(E'') + \chern(E')$. Furthermore,
    \[\chern: K(X) \longrightarrow H^*(X;\Q)\]
    is a ring homomorphism, i.e., $\chern(E\otimes F)=\chern(E)\chern(F)$.
\end{prop}

\begin{proof}
    By the splitting principle, assume $E',E,E''$ filter by line bundles, so their Chern polynomials split, and we have respective (multi)sets of Chern roots $A', A, A''$. The Whitney sum formula says exactly $A = A' \cup A''$, so
    \[\chern(E) = \sum_{a\in A} e^{at} = \sum_{a'\in A'} e^{a't} + \sum_{a''\in A''} e^{a''t} = \chern(E') + \chern(E'').\]
    So, $\chern$ is additive on exact sequences, to it descends to a group homomorphism $K(X) \to H^*(X;\Q)$.

    We now show that $\chern$ is in fact a ring homomorphism. So, we compute the Chern character of a tensor product. Let $E,F$ be vector bundles with Chern roots $a_i,b_j$ respectively. Then,
    \begin{align*}
        \chern(E\otimes F) &= \sum_{i,j} e^{(a_i+b_j)t}\\
        &= \sum_{i,j} e^{a_it} e^{b_j t}\\
        &= \left(\sum_i e^{a_i t}\right)\left(\sum_j e^{b_j t}\right)\\
        &= \chern(E) \chern(F).
    \end{align*}
\end{proof}

\begin{eg}\label{skyscraper-on-p2}
    We compute $\chern$ of the skyscraper sheaf $\C_P$ for a point $P$ on $\P^2$. We have a Koszul complex
    \[0 \to \cO_{\P^2}(-2) \to \cO_{\P^2}(-1)^{\oplus 2} \to \cO_{\P^2} \to \C_P \to 0,\]
    so by additivity of $\chern$, we get
    \begin{align*}
        \chern(\C_P) &= \chern(\cO_{\P^2}) - \chern(\cO_{\P^2}(-1)^{\oplus 2}) + \chern(\cO_{\P^2}(-2))\\
        &= e^{0} - 2e^{-Ht} + e^{-2Ht}\\
        &= 1 - 2(1-Ht+\frac{1}{2}H^2 t^2) + (1-2Ht + \frac{1}{2}(2H)^2t^2)\\
        &= 0 + 0t + H^2 t^2\\
        &= t^2
    \end{align*}
    as an element of $H^*(\P^2;\Q) = \Q[H]/(H^{2+1})$. Sometimes, we will identify $H^*(\P^2;\Q)\cong \Q^3$ as a $\Q$-module, and say that $\chern(\C_P)$ is the vector $(0,0,1)$.
\end{eg}



\section{Grothendieck Riemann Roch and Applications}

The Chern character is defined on the Grothendieck group $X$, and, for example, on a projective variety, all coherent sheaves have finite resolutions by vector bundles (see Hilbert syzygy theorem). This will identify the Grothendieck group of vector bundles on $X$ with the Grothendieck group of coherent sheaves.

Our knowledge on computing Chern characters is pretty limited so far. We at least know how to compute the Chern character of a vector bundle that has a filtration by line bundles, as in the splitting principle. We also have example \ref{skyscraper-on-p2}, where we computed the skyscraper sheaf of a particularly simple pushforward. We will see how generalize this example to pushforwards of structure sheaves under closed embeddings via a generalization of Riemann-Roach.

First, let's understand how the classical Riemann-Roch statement is actually a statement about Chern characters, and then we will state Grothendieck Riemann-Roch.

\begin{thm}[Riemann-Roch]
    Let $D$ be a divisor on a curve $C$ of genus $g$. Then,
    \[\chi(\cO_C(D)) = \deg D + 1 - g,\]
    where $\chi = \sum (-1)^i h^i$ is the Euler characteristic.
\end{thm}

We will cast all the terms in new light.

First, we reinterpret $\chi(\cO_C(D))$. Consider $\pi: C\to \Spec \C$, then sheaf cohomology $H^i$ can be identified with higher direct images $R^i\pi_*$. If we think of $\cF=\cO_C(D)$ as an object in the derived category, $D^b(C)$, then $R\pi_* \cF$ lives in $D^b(\Spec \C)$ is a bounded complex of vector spaces. Then, we can pass to the Grothendieck group by taking an alternating formal sum
\begin{align*}
    D^b(\Spec \C) &\longrightarrow K(\Spec \C)\\
    (0\to V^{-n} \to \cdots \to V^n \to 0) &\longmapsto \sum_{i=-n}^n (-1)^i [V^i],
\end{align*}
and finally we get the Euler characteristic as the rank homomorphism $K(\Spec \C) \to \Z$, which is same thing as the (zeroth) Chern character, where $H^*(\Spec\C) = \Z$. So, in total, we should identify
\[\chi(\cF) = \chern(R\pi_* \cF).\]

The terms $\deg D$ and $1-g$ will then be gotten by taking first Chern classes of $D$ and $\frac{-1}{2} K_X$, which will pop out of the generalization of Riemann-Roch. To get the genus (so the $K_X$ term), we will define the following variant of Chern classes.


\begin{defn}
    For a vector bundle $E$ over $X$, there is an element $\todd(E) \in H^*(X;\Q)$ called the \textbf{Todd class}, defined as follows. By the splitting principle, assume $E$ has Chern roots $a_i$, and set
    \[\todd(E) = \prod_{i=1}^r \frac{a_i t}{1-e^{-a_i t}},\]
    where
    \[\frac{x}{1-e^{-x}} = 1 + \frac{1}{2}x + \frac{1}{12}x^2 - \frac{1}{720}x^4 + \cdots.\]
    One can work out that
    \[\todd = 1 + \frac{1}{2}c_1 t + \frac{1}{12}(c_1^2+c_2)t^2 + \frac{1}{24}c_1c_2 t^3 -\frac{1}{720}(c_1^4-4c_1^2c_2-3c_2^2-c_1c_3+c_4)t^4 + \cdots\]
\end{defn}

\begin{rmk}
    A todd class $\todd(\cF)$ is a unit in $H^*(X;\Q)$, since it is a sum of a unit $1$ and a nilpotent $\todd(\cF)-1$ (since it has a factor of $t$, and $t^{\dim X+1}=0$). So, it makes sense to divide by Todd classes.
\end{rmk}

\begin{prop}
    Todd classes are multiplicative on exact sequences.
\end{prop}

\begin{thm}[Grothendieck-Riemann-Roch]
    Let $X,Y$ be smooth quasi-projective varieties, $f:X\to Y$ a proper morphism. Then,
    \[\chern(Rf_* \cF^\bullet) \todd(\cT_Y) = f_*(\chern(\cF^\bullet)\todd(\cT_X)).\]
\end{thm}

If we specialize to $\pi: X\to \Spec k$ the structure morphism for a smooth projective variety over $k$, we recover Hirzebruch-Riemann-Roch.
\begin{eg}[Hirzebruch-Riemann-Roch]
    We already saw that
    \[\chern(R\pi_* \cF) = \chi(\cF) \in H^*(\Spec \C; \Q) \cong \Q,\]
    and on the point $\Spec \C$, $\cT_{\Spec \C}=\cO$, so $\todd(\cT_{\Spec \C}) = \todd(\cO) = 1$. So, the right hand side of GRR becomes just $\chi(\cF)$.

    Now, for the right hand side, we have to compute cohomology pushforward. The class $\chern(\cF)\todd(\cT_X) \in H^*(X;\Q)$ has a Poincare dual $\alpha \in H_*(X;\Q)$. For the pushforward $f_*: H_*(X;\Q) \to H_*(\Spec \C; \Q)$, only the bottom homology $H_0(\Spec \C; \Q) \cong \Q$ exists, which Poincare duality sends isomorphically to $H^1(\Spec \C; \Q)$.
    
    In total, we only need the top cohomology component of $\chern(\cF)\todd(\cT_X)$, which is sent by Poincare duality to
    \[\left(\int_X \chern(\cF)\todd(\cT_X)\right)[\mathrm{pt}],\]
    (interpreting $\chern(\cF)\todd(\cT_X)$ as a sum of forms), which is eventually sent to just
    \[\int_X \chern(\cF)\todd(\cT_X)\]
    once we identify $H^0(\Spec \C;\Q) \cong \Q$ (the isomorphism sending the constant 1-function $1: \Spec \C\to \Q$ to $1\in \Q$).
    
    So, GRR specializes to
    \[\chi(\cF) = \int \chern(\cF) \todd(\cT_X).\]
\end{eg}

Now, let's specialize even farther to get Riemann-Roch for curves.

\begin{eg}[Proof of Riemann-Roch]
    Let $D$ be a divisor on a curve $C$, and $\pi:C\to \Spec\C$ the structure map. We have already established that $\chern(R\pi_* \cO_C(D)) = \chi(\cO_C(D))$, and $\todd(\cT_{\Spec \C})=\todd(0)=1$. So, the left hand side is indeed $\chi(\cO_C(D))$.

    Now, we compute the right hand side. The line bundle $\cO_C(D)$ has Chern root $D$, so its Chern character is just
    \[\chern(\cO_C(D)) = 1+Dt.\]
    Next,
    \begin{align*}
        \todd(\cT_X) &= \todd(\cO(-K_C))\\
        &= 1 -\frac{1}{2}K_C t,
    \end{align*}
    so
    \[\chern(\cO_C(D))\todd(\cT_X) = 1 + (D-\frac{1}{2}K_C)t.\]
    Now, we pushforward. We do this by taking Poincare duals, homology pushforward, and Poincare duals again.
    \[(1,0,D-\frac{1}{2}K_C) \mapsto (\deg(D-\frac{1}{2}K_C),0,1) \mapsto \deg(D-\frac{1}{2}K_C) \mapsto \deg(D-\frac{1}{2}K_C) = \deg D - (g-1),\]
    so equating both sides we exactly get
    \[\chi(\cO_C(D)) = \deg D + 1 - g.\]
\end{eg}

\begin{eg}[Pushforwards Along Closed Embeddings]
    Let $f:X\to Y$ be a closed embedding. We will calculate the Chern character of $f_*\cO_X$. We go through and reinterpret each term in GRR.
    
    First, we work on the left hand side. Since $f$ is a closed embedding, it is affine, so there is no higher direct images, so $\chern(Rf_* \cO_X) = \chern(f_* \cO_X)$.

    To handle the right hand side, we first use the normal bundle sequence
    \[0\to \cT_X \to f^*\cT_Y \to \cN_{X/Y} \to 0\]
    so
    \[\todd(\cT_X) = \frac{\todd(f^* \cT_Y)}{\todd(\cN_{X/Y})} = \frac{f^*(\todd(\cT_Y))}{\todd(\cN_{X/Y})}.\]
    Now, we use a projection formula for singular cohomology classes (at least for those represented by algebraic cycles??)
    \[f_*(x \cdot f^* y) = f_*(x) \cdot y,\]
    so we can write
    \begin{align*}
        f_*(\chern(\cO_X)\todd(\cT_X)) &= f_*\left(\frac{\chern(\cO_X)}{\todd(\cN_{X/Y})}\cdot f^*(\todd(\cT_Y))\right)\\
        &= f_*\left(\frac{\chern(\cO_X)}{\todd(\cN_{X/Y})}\right) \todd(\cT_Y).
    \end{align*}
    Additionally, $\chern(\cO_X)=e^{0t}=1$, so this term is just
    \[f_*\left(\frac{1}{\todd(\cN_{X/Y})}\right) \todd(\cT_Y).\]
    Thus, canceling out the $\todd(\cT_Y)$ on both sides (since it is a unit) we get
    \[\chern(f_* \cO_X) = f_*\left(\frac{1}{\todd(\cN_{X/Y})}\right).\]

    For example, when $X$ is a smooth point and $Y$ is projective of dimension $n$, then $\cN_{X/Y} = \cO_X^{\oplus n}$, and so
    \[\todd(\cN_{X/Y}) = 1,\]
    which is pushed forward (i.e., do Poincare duality, homology pushforward, then Poincare duality) to
    \[t^n \in H^*(Y;\Q),\]
    giving us the Chern character of $f_*\cO_X$.
\end{eg}

\nocite{*} % Include all references in refs.bib regardless of whether they are referenced or not
\bibliographystyle{plain} % We choose the "plain" reference style
\bibliography{2024fall/intersection_theory_refs} % Entries are in the refs.bib file

\end{document}
